\documentclass[12pt,a4paper]{article}
\usepackage[utf8]{inputenc}
\usepackage{graphicx}
\usepackage{listings}
\usepackage{xcolor}
\usepackage{hyperref}
\usepackage{geometry}
\usepackage{booktabs}
\usepackage{enumitem}
\usepackage{fancyhdr}
\usepackage{titlesec}
\usepackage{amsmath}
\usepackage{float}

\geometry{margin=2.5cm}

% Colors
\definecolor{codegreen}{rgb}{0,0.6,0}
\definecolor{codegray}{rgb}{0.5,0.5,0.5}
\definecolor{codepurple}{rgb}{0.58,0,0.82}
\definecolor{backcolour}{rgb}{0.95,0.95,0.92}
\definecolor{headerblue}{RGB}{44,62,80}
\definecolor{accentgreen}{RGB}{39,174,96}
\definecolor{accentblue}{RGB}{52,152,219}
\definecolor{accentred}{RGB}{231,76,60}

% Code style
\lstdefinestyle{mystyle}{
    backgroundcolor=\color{backcolour},
    commentstyle=\color{codegreen},
    keywordstyle=\color{blue},
    numberstyle=\tiny\color{codegray},
    stringstyle=\color{codepurple},
    basicstyle=\ttfamily\footnotesize,
    breakatwhitespace=false,
    breaklines=true,
    captionpos=b,
    keepspaces=true,
    numbers=left,
    numbersep=5pt,
    showspaces=false,
    showstringspaces=false,
    showtabs=false,
    tabsize=2,
    frame=single
}
\lstset{style=mystyle}

% JavaScript language definition
\lstdefinelanguage{JavaScript}{
  keywords={const, let, var, function, return, if, else, for, while, do, break, continue, new, try, catch, throw, async, await, class, extends, super, this, typeof, instanceof, switch, case, default, null, undefined, true, false},
  morecomment=[l]{//},
  morecomment=[s]{/*}{*/},
  morestring=[b]',
  morestring=[b]",
  morestring=[b]`,
}

% JSON language definition
\lstdefinelanguage{JSON}{
  keywords={true, false, null},
  morecomment=[l]{//},
  morecomment=[s]{/*}{*/},
  morestring=[b]",
}

% CSS language definition
\lstdefinelanguage{CSS}{
  keywords={color,background,background-color,border,border-radius,display,position,fixed,flex,justify-content,align-items,padding,margin,width,height,font-size,font-weight,border-collapse,overflow-x,transition,z-index,min-height,max-width,cursor,text-align,text-transform,letter-spacing},
  morecomment=[l]{//},
  morecomment=[s]{/*}{*/},
  morestring=[b]",
  morestring=[b]',
}

% Header/Footer
\pagestyle{fancy}
\fancyhf{}
\fancyhead[L]{\small Online Store Dashboard}
\fancyhead[R]{\small Student Guide}
\fancyfoot[C]{\thepage}

% Title formatting
\titleformat{\section}{\Large\bfseries\color{headerblue}}{\thesection}{1em}{}
\titleformat{\subsection}{\large\bfseries\color{accentblue}}{\thesubsection}{1em}{}

\begin{document}

% ==================== TITLE PAGE ====================
\begin{titlepage}
    \centering
    \vspace*{2cm}
    
    {\Huge\bfseries Online Store Dashboard\par}
    \vspace{1cm}
    {\Large Complete Student Guide\par}
    \vspace{0.5cm}
    {\large Express.js + Supabase + HTML/CSS/JavaScript\par}
    
    \vspace{2cm}
    
    \begin{center}
        \fbox{
            \begin{minipage}{0.8\textwidth}
                \centering
                \vspace{0.5cm}
                {\large\textbf{Project Overview}\par}
                \vspace{0.3cm}
                Is project mein humne ek complete Online Store Dashboard banaya hai\\
                jisme products ko add, edit, delete aur view kar sakte hain.\\
                Pehle local server (Express) use kiya, phir Supabase (cloud) par migrate kiya.
                \vspace{0.5cm}
            \end{minipage}
        }
    \end{center}
    
    \vspace{2cm}
    {\large\textbf{Technologies Used}\par}
    \vspace{0.5cm}
    \begin{tabular}{ll}
        \textbf{Backend:} & Node.js, Express.js, CORS \\
        \textbf{Database:} & JSON file (local), Supabase (cloud) \\
        \textbf{Frontend:} & HTML5, CSS3, JavaScript (ES6+) \\
        \textbf{API:} & REST API, Supabase JavaScript Client \\
    \end{tabular}
    
    \vfill
    {\large\date{\today}\par}
\end{titlepage}

% ==================== TABLE OF CONTENTS ====================
\newpage
\tableofcontents
\newpage

% ==================== SECTION 1 ====================
\section{Project Structure}
\label{sec:structure}

\subsection{Folder Structure}
Sab se pehle humne project ka folder structure banaya:

\begin{lstlisting}[language=bash, caption={Project Folder Structure}]
online-store-dashboard/
|-- backend/
|   |-- package.json
|   |-- package-lock.json
|   |-- products.json
|   `-- server.js
`-- frontend/
    |-- index.html
    |-- style.css
    `-- script.js
\end{lstlisting}

\begin{itemize}[leftmargin=2em]
    \item \textbf{backend/} - Server-side code, data storage, API routes
    \item \textbf{frontend/} - User interface, styling, client-side logic
    \item \textbf{products.json} - Sample products ka data (3 products)
\end{itemize}

\subsection{Products Data}
\texttt{products.json} mein 3 sample products hain:

\begin{lstlisting}[language=json, caption={products.json - Sample Data}]
[
  {
    "id": 1,
    "name": "Wireless Mouse",
    "price": 1500,
    "stock": 25,
    "category": "Electronics"
  },
  {
    "id": 2,
    "name": "Notebook",
    "price": 200,
    "stock": 100,
    "category": "Stationery"
  },
  {
    "id": 3,
    "name": "Water Bottle",
    "price": 500,
    "stock": 50,
    "category": "Accessories"
  }
]
\end{lstlisting}

% ==================== SECTION 2 ====================
\section{Express Server Setup}
\label{sec:express}

\subsection{Installation}
Backend folder mein Express install karne ke liye:

\begin{lstlisting}[language=bash, caption={Express Installation}]
cd backend
npm init -y
npm install express
npm install cors
\end{lstlisting}

\subsection{Server Code}
\texttt{server.js} - Basic Express server:

\begin{lstlisting}[language=JavaScript, caption={server.js - Express Server}]
const express = require('express');
const cors = require('cors');
const fs = require('fs');
const path = require('path');

const app = express();

// Middleware
app.use(cors());
app.use(express.json());

// products.json file ka path
const productsFile = path.join(__dirname, 'products.json');

// Helper functions
function readProducts() {
  const data = fs.readFileSync(productsFile, 'utf8');
  return JSON.parse(data);
}

function writeProducts(products) {
  fs.writeFileSync(productsFile, JSON.stringify(products, null, 2));
}

// Root route
app.get('/', (req, res) => {
  res.send('Welcome to Online Store API!');
});

app.listen(3000, () => {
  console.log('Server chal raha hai: http://localhost:3000');
});
\end{lstlisting}

\subsection{Server Chalana}
\begin{lstlisting}[language=bash, caption={Server Start}]
node server.js
\end{lstlisting}

Output: \texttt{Server chal raha hai: http://localhost:3000}

% ==================== SECTION 3 ====================
\section{CRUD API (Backend)}
\label{sec:crud}

\subsection{REST API Routes}
CRUD ka matlab hai: \textbf{C}reate, \textbf{R}ead, \textbf{U}pdate, \textbf{D}elete.

\begin{table}[H]
\centering
\caption{CRUD API Routes}
\begin{tabular}{@{}llll@{}}
\toprule
\textbf{Method} & \textbf{Route} & \textbf{Kaam} & \textbf{Status} \\
\midrule
GET & /api/products & Saare products & 200 \\
GET & /api/products/:id & Ek product & 200/404 \\
POST & /api/products & Naya product add & 201 \\
PUT & /api/products/:id & Product update & 200/404 \\
DELETE & /api/products/:id & Product delete & 200/404 \\
\bottomrule
\end{tabular}
\end{table}

\subsection{CRUD Routes Code}

\begin{lstlisting}[language=JavaScript, caption={GET - Saare Products}]
// GET /api/products -> saare products return karta hai (200)
app.get('/api/products', (req, res) => {
  const products = readProducts();
  res.status(200).json(products);
});
\end{lstlisting}

\begin{lstlisting}[language=JavaScript, caption={GET - Ek Product by ID}]
// GET /api/products/:id -> ek product return karta hai
app.get('/api/products/:id', (req, res) => {
  const products = readProducts();
  const id = parseInt(req.params.id);
  const product = products.find(p => p.id === id);

  if (!product) {
    return res.status(404).json({ message: 'Product nahi mila' });
  }

  res.status(200).json(product);
});
\end{lstlisting}

\begin{lstlisting}[language=JavaScript, caption={POST - Naya Product Add}]
// POST /api/products -> naya product add karta hai (201)
app.post('/api/products', (req, res) => {
  const products = readProducts();
  const { name, price, stock, category } = req.body;

  const newId = products.length > 0 
    ? Math.max(...products.map(p => p.id)) + 1 
    : 1;

  const newProduct = { id: newId, name, price, stock, category };

  products.push(newProduct);
  writeProducts(products);

  res.status(201).json(newProduct);
});
\end{lstlisting}

\begin{lstlisting}[language=JavaScript, caption={PUT - Product Update}]
// PUT /api/products/:id -> product update karta hai
app.put('/api/products/:id', (req, res) => {
  const products = readProducts();
  const id = parseInt(req.params.id);
  const index = products.findIndex(p => p.id === id);

  if (index === -1) {
    return res.status(404).json({ message: 'Product nahi mila' });
  }

  const { name, price, stock, category } = req.body;

  if (name !== undefined) products[index].name = name;
  if (price !== undefined) products[index].price = price;
  if (stock !== undefined) products[index].stock = stock;
  if (category !== undefined) products[index].category = category;

  writeProducts(products);
  res.status(200).json(products[index]);
});
\end{lstlisting}

\begin{lstlisting}[language=JavaScript, caption={DELETE - Product Delete}]
// DELETE /api/products/:id -> product delete karta hai
app.delete('/api/products/:id', (req, res) => {
  const products = readProducts();
  const id = parseInt(req.params.id);
  const index = products.findIndex(p => p.id === id);

  if (index === -1) {
    return res.status(404).json({ message: 'Product nahi mila' });
  }

  const deletedProduct = products.splice(index, 1)[0];
  writeProducts(products);

  res.status(200).json({ message: 'Product delete ho gaya', deletedProduct });
});
\end{lstlisting}

% ==================== SECTION 4 ====================
\section{Frontend Design (HTML + CSS)}
\label{sec:frontend}

\subsection{HTML Structure}
\texttt{index.html} - Dashboard ka main structure:

\begin{lstlisting}[language=HTML, caption={index.html - Main Structure}]
<!DOCTYPE html>
<html lang="en">
<head>
  <meta charset="UTF-8">
  <meta name="viewport" content="width=device-width, initial-scale=1.0">
  <title>Online Store Dashboard</title>
  <link rel="stylesheet" href="style.css">
</head>
<body>
  <!-- Toast Notifications Container -->
  <div id="toastContainer" class="toast-container"></div>

  <!-- Header -->
  <header class="header">
    <h1>Online Store Dashboard</h1>
  </header>

  <!-- Main Container -->
  <main class="container">
    <!-- Toolbar -->
    <div class="toolbar">
      <button id="addProductBtn" class="btn btn-green">+ Add Product</button>
      <input type="text" id="searchInput" class="search-input" placeholder="Search products...">
    </div>

    <!-- Products Table -->
    <div class="table-wrapper">
      <table class="products-table" id="productsTable">
        <thead>
          <tr>
            <th>ID</th>
            <th>Name</th>
            <th>Price</th>
            <th>Stock</th>
            <th>Category</th>
            <th>Actions</th>
          </tr>
        </thead>
        <tbody id="productBody">
          <!-- Products yahan JavaScript se add honge -->
        </tbody>
      </table>
    </div>
  </main>

  <!-- Modal (Add/Edit Product) -->
  <div id="productModal" class="modal-overlay hidden">
    <div class="modal">
      <h2 id="modalTitle">Add Product</h2>
      <form id="productForm">
        <div class="form-group">
          <label for="productName">Name</label>
          <input type="text" id="productName" name="name" required>
        </div>
        <div class="form-group">
          <label for="productPrice">Price</label>
          <input type="number" id="productPrice" name="price" required>
        </div>
        <div class="form-group">
          <label for="productStock">Stock</label>
          <input type="number" id="productStock" name="stock" required>
        </div>
        <div class="form-group">
          <label for="productCategory">Category</label>
          <input type="text" id="productCategory" name="category" required>
        </div>
        <div class="modal-actions">
          <button type="submit" class="btn btn-green">Save</button>
          <button type="button" id="cancelBtn" class="btn btn-gray">Cancel</button>
        </div>
      </form>
    </div>
  </div>

  <!-- Supabase CDN -->
  <script src="https://cdn.jsdelivr.net/npm/@supabase/supabase-js@2"></script>
  <script src="script.js"></script>
</body>
</html>
\end{lstlisting}

\subsection{CSS Styling}
\texttt{style.css} - Modern clean design:

\begin{itemize}[leftmargin=2em]
    \item \textbf{Colors:} Light gray background (\#f0f2f5), white cards
    \item \textbf{Buttons:} Green (Add), Blue (Edit), Red (Delete)
    \item \textbf{Table:} Dark header, hover effect
    \item \textbf{Modal:} Dark overlay + centered white form
    \item \textbf{Responsive:} Mobile par toolbar stack, table scroll
\end{itemize}

\begin{lstlisting}[language=CSS, caption={style.css - Key Styles}]
/* Toast Notifications */
.toast-container {
  position: fixed;
  top: 20px;
  right: 20px;
  z-index: 2000;
}

.toast-success { background-color: #27ae60; }
.toast-error { background-color: #e74c3c; }

/* Buttons */
.btn-green { background-color: #27ae60; color: #ffffff; }
.btn-edit { background-color: #3498db; color: #ffffff; }
.btn-delete { background-color: #e74c3c; color: #ffffff; }

/* Table */
.products-table thead { background-color: #2c3e50; color: #ffffff; }
.products-table tbody tr:hover { background-color: #f5f7fa; }

/* Modal */
.modal-overlay {
  position: fixed;
  background-color: rgba(0, 0, 0, 0.6);
  display: flex;
  justify-content: center;
  align-items: center;
}

/* Responsive */
@media (max-width: 768px) {
  .toolbar { flex-direction: column; }
  .btn, .btn-edit, .btn-delete { min-height: 40px; }
}
\end{lstlisting}

% ==================== SECTION 5 ====================
\section{Frontend JavaScript (CRUD Logic)}
\label{sec:js}

\subsection{Supabase Client Initialization}
\texttt{script.js} - Supabase client setup:

\begin{lstlisting}[language=JavaScript, caption={Supabase Client Setup}]
// ===== Supabase Client Initialization =====

const supabaseUrl = 'https://dntjgvacurwpcyavwdkc.supabase.co';
const supabaseKey = 'sb_publishable_-_6VQDTsincVrv5FwvCZow_g1F7dFHA';
const supabase = supabase.createClient(supabaseUrl, supabaseKey);

// DOMContentLoaded par loadProducts() call karo
document.addEventListener('DOMContentLoaded', () => {
  loadProducts();
});
\end{lstlisting}

\subsection{Toast Notifications}
\begin{lstlisting}[language=JavaScript, caption={Toast Notification Function}]
// showToast(): success/error message top-right me dikhata hai
function showToast(message, type = 'success') {
  const container = document.getElementById('toastContainer');

  const toast = document.createElement('div');
  toast.className = `toast toast-${type}`;
  toast.textContent = message;

  container.appendChild(toast);

  // 3 second baad auto-remove
  setTimeout(() => {
    toast.classList.add('hide');
    setTimeout(() => {
      toast.remove();
    }, 300);
  }, 3000);
}
\end{lstlisting}

\subsection{Read Products (GET)}
\begin{lstlisting}[language=JavaScript, caption={loadProducts() - Supabase Select}]
// loadProducts(): Supabase se products fetch karta hai
async function loadProducts() {
  const tbody = document.getElementById('productBody');

  // Loading state dikhao
  tbody.innerHTML = '<tr><td colspan="6" style="text-align:center; color:#888;">Loading...</td></tr>';

  try {
    const { data: products, error } = await supabase
      .from('products')
      .select('*');

    if (error) throw error;

    renderProducts(products);
  } catch (error) {
    console.error('Products load karne me error:', error);
    tbody.innerHTML = '<tr><td colspan="6" style="text-align:center; color:#e74c3c;">Products load karne me error aaya. Server check karein.</td></tr>';
    showToast('Error: Server se connect nahi ho paya', 'error');
  }
}
\end{lstlisting}

\subsection{Render Products}
\begin{lstlisting}[language=JavaScript, caption={renderProducts() - Table Rows}]
// renderProducts(): products ko table rows me render karta hai
function renderProducts(products) {
  const tbody = document.getElementById('productBody');

  // Agar products khali hain to message dikhao
  if (products.length === 0) {
    tbody.innerHTML = '<tr><td colspan="6" style="text-align:center; color:#888;">Koi product nahi hai</td></tr>';
    return;
  }

  // Har product ki ek table row banayein
  const rows = products.map(product => `
    <tr>
      <td>${product.id}</td>
      <td>${product.name}</td>
      <td>Rs. ${product.price}</td>
      <td>${product.stock}</td>
      <td>${product.category}</td>
      <td>
        <button class="btn-edit" data-id="${product.id}" 
                data-name="${product.name}" data-price="${product.price}" 
                data-stock="${product.stock}" data-category="${product.category}">Edit</button>
        <button class="btn-delete" data-id="${product.id}">Delete</button>
      </td>
    </tr>
  `).join('');

  tbody.innerHTML = rows;
}
\end{lstlisting}

\subsection{Add Product (POST)}
\begin{lstlisting}[language=JavaScript, caption={Form Submit - Add/Edit}]
document.getElementById('productForm').addEventListener('submit', async (e) => {
  e.preventDefault();

  const name = document.getElementById('productName').value.trim();
  const price = parseFloat(document.getElementById('productPrice').value);
  const stock = parseInt(document.getElementById('productStock').value);
  const category = document.getElementById('productCategory').value.trim();

  // Form validation
  if (!name || !category) {
    showToast('Error: Name aur Category zaroori hain!', 'error');
    return;
  }

  if (isNaN(price) || price < 0) {
    showToast('Error: Price negative nahi ho sakta!', 'error');
    return;
  }

  if (isNaN(stock) || stock < 0) {
    showToast('Error: Stock negative nahi ho sakta!', 'error');
    return;
  }

  const productData = { name, price, stock, category };

  try {
    if (editingId !== null) {
      // Edit mode: Supabase update karo
      const { error } = await supabase
        .from('products')
        .update(productData)
        .eq('id', editingId);

      if (error) {
        showToast(`Error: ${error.message || 'Update me error aaya'}`, 'error');
        return;
      }
    } else {
      // Add mode: Supabase insert karo
      const { error } = await supabase
        .from('products')
        .insert([productData]);

      if (error) {
        showToast(`Error: ${error.message || 'Add karne me error aaya'}`, 'error');
        return;
      }
    }

    // Success par: modal band, loadProducts() call
    document.getElementById('productModal').classList.add('hidden');
    document.getElementById('productForm').reset();
    editingId = null;
    loadProducts();
    showToast('Product saved ho gaya!', 'success');
  } catch (error) {
    console.error('Error:', error);
    showToast('Error: Server se connect nahi ho paya', 'error');
  }
});
\end{lstlisting}

\subsection{Delete Product (DELETE)}
\begin{lstlisting}[language=JavaScript, caption={deleteProduct() - Supabase Delete}]
// deleteProduct(): Supabase se product delete karta hai
async function deleteProduct(id) {
  try {
    const { error } = await supabase
      .from('products')
      .delete()
      .eq('id', id);

    if (error) {
      showToast(`Error: ${error.message || 'Delete me error aaya'}`, 'error');
      return;
    }

    // Delete ke baad loadProducts() call karo
    loadProducts();
    showToast('Product delete ho gaya!', 'success');
  } catch (error) {
    console.error('Error:', error);
    showToast('Error: Server se connect nahi ho paya', 'error');
  }
}
\end{lstlisting}

% ==================== SECTION 6 ====================
\section{Supabase Migration}
\label{sec:supabase}

\subsection{Fetch vs Supabase Comparison}
\begin{table}[H]
\centering
\caption{Fetch API vs Supabase Methods}
\begin{tabular}{@{}lll@{}}
\toprule
\textbf{Operation} & \textbf{Fetch (Local)} & \textbf{Supabase} \\
\midrule
Read & \texttt{fetch('/api/products')} & \texttt{supabase.from('products').select('*')} \\
Create & \texttt{fetch('/api/products', \{method:'POST'\})} & \texttt{supabase.from('products').insert([data])} \\
Update & \texttt{fetch('/api/products/:id', \{method:'PUT'\})} & \texttt{supabase.from('products').update(data).eq('id', id)} \\
Delete & \texttt{fetch('/api/products/:id', \{method:'DELETE'\})} & \texttt{supabase.from('products').delete().eq('id', id)} \\
\bottomrule
\end{tabular}
\end{table}

\subsection{Supabase Setup Steps}
\begin{enumerate}[leftmargin=2em]
    \item Supabase account banayein
    \item Naya project create karein
    \item \texttt{products} table banayein (id, name, price, stock, category)
    \item Project URL aur anon key copy karein
    \item CDN script HTML mein add karein
    \item Client initialize karein
\end{enumerate}

% ==================== SECTION 7 ====================
\section{Testing}
\label{sec:testing}

\subsection{API Testing (Thunder Client / Postman)}
\begin{table}[H]
\centering
\caption{API Test Results}
\begin{tabular}{@{}llll@{}}
\toprule
\textbf{Method} & \textbf{URL} & \textbf{Body} & \textbf{Result} \\
\midrule
GET & /api/products & - & 200 - Saare products \\
GET & /api/products/1 & - & 200 - Ek product \\
GET & /api/products/999 & - & 404 - Product nahi mila \\
POST & /api/products & \{name, price, stock, category\} & 201 - Naya product \\
PUT & /api/products/4 & \{price: 2800\} & 200 - Update hua \\
DELETE & /api/products/4 & - & 200 - Delete hua \\
\bottomrule
\end{tabular}
\end{table}

\subsection{Browser Testing}
\begin{enumerate}[leftmargin=2em]
    \item Server chalayein: \texttt{node server.js}
    \item Browser mein \texttt{index.html} kholen
    \item Products table mein dikhne chahiye
    \item Add/Edit/Delete test karein
    \item Mobile view (DevTools) mein responsive check karein
\end{enumerate}

% ==================== SECTION 8 ====================
\section{Seekhne Ke Liye Key Concepts}
\label{sec:learn}

\begin{enumerate}[leftmargin=2em]
    \item \textbf{Project Structure} - Backend/frontend alag rakhna
    \item \textbf{REST API} - CRUD operations ka concept
    \item \textbf{Express.js} - Node.js server banana
    \item \textbf{Fetch API} - Frontend se backend data lena
    \item \textbf{CORS} - Cross-origin requests
    \item \textbf{DOM Manipulation} - Table rows, modal, event listeners
    \item \textbf{Async/Await} - Asynchronous JavaScript
    \item \textbf{Form Validation} - User input check karna
    \item \textbf{Responsive Design} - Mobile-friendly UI
    \item \textbf{Toast Notifications} - User feedback
    \item \textbf{Supabase} - Cloud database (Backend-as-a-Service)
    \item \textbf{Migration} - Local se cloud par shift karna
\end{enumerate}

% ==================== SECTION 9 ====================
\section{Next Steps / Aage Kya Kar Sakte Hain}
\label{sec:next}

\begin{itemize}[leftmargin=2em]
    \item Search functionality implement karna
    \item Supabase authentication add karna
    \item Product images add karna
    \item Sorting/filtering features
    \item Deploy karna (Netlify/Vercel for frontend)
    \item Pagination add karna
    \item Real-time updates (Supabase Realtime)
\end{itemize}

% ==================== APPENDIX ====================
\appendix
\section{Complete File Listings}
\label{sec:appendix}

\subsection{server.js (Complete)}
\begin{lstlisting}[language=JavaScript, caption={server.js - Full Code}]
const express = require('express');
const cors = require('cors');
const fs = require('fs');
const path = require('path');

const app = express();

// Middleware
app.use(cors());
app.use(express.json());

const productsFile = path.join(__dirname, 'products.json');

function readProducts() {
  const data = fs.readFileSync(productsFile, 'utf8');
  return JSON.parse(data);
}

function writeProducts(products) {
  fs.writeFileSync(productsFile, JSON.stringify(products, null, 2));
}

// Root route
app.get('/', (req, res) => {
  res.send('Welcome to Online Store API!');
});

// GET /api/products
app.get('/api/products', (req, res) => {
  const products = readProducts();
  res.status(200).json(products);
});

// GET /api/products/:id
app.get('/api/products/:id', (req, res) => {
  const products = readProducts();
  const id = parseInt(req.params.id);
  const product = products.find(p => p.id === id);

  if (!product) {
    return res.status(404).json({ message: 'Product nahi mila' });
  }

  res.status(200).json(product);
});

// POST /api/products
app.post('/api/products', (req, res) => {
  const products = readProducts();
  const { name, price, stock, category } = req.body;

  const newId = products.length > 0 
    ? Math.max(...products.map(p => p.id)) + 1 
    : 1;

  const newProduct = { id: newId, name, price, stock, category };

  products.push(newProduct);
  writeProducts(products);

  res.status(201).json(newProduct);
});

// PUT /api/products/:id
app.put('/api/products/:id', (req, res) => {
  const products = readProducts();
  const id = parseInt(req.params.id);
  const index = products.findIndex(p => p.id === id);

  if (index === -1) {
    return res.status(404).json({ message: 'Product nahi mila' });
  }

  const { name, price, stock, category } = req.body;

  if (name !== undefined) products[index].name = name;
  if (price !== undefined) products[index].price = price;
  if (stock !== undefined) products[index].stock = stock;
  if (category !== undefined) products[index].category = category;

  writeProducts(products);
  res.status(200).json(products[index]);
});

// DELETE /api/products/:id
app.delete('/api/products/:id', (req, res) => {
  const products = readProducts();
  const id = parseInt(req.params.id);
  const index = products.findIndex(p => p.id === id);

  if (index === -1) {
    return res.status(404).json({ message: 'Product nahi mila' });
  }

  const deletedProduct = products.splice(index, 1)[0];
  writeProducts(products);

  res.status(200).json({ message: 'Product delete ho gaya', deletedProduct });
});

app.listen(3000, () => {
  console.log('Server chal raha hai: http://localhost:3000');
});
\end{lstlisting}

\subsection{script.js (Complete)}
\begin{lstlisting}[language=JavaScript, caption={script.js - Full Code}]
// ===== Supabase Client Initialization =====

const supabaseUrl = 'https://dntjgvacurwpcyavwdkc.supabase.co';
const supabaseKey = 'sb_publishable_-_6VQDTsincVrv5FwvCZow_g1F7dFHA';
const supabase = supabase.createClient(supabaseUrl, supabaseKey);

// DOMContentLoaded par loadProducts() call karo
document.addEventListener('DOMContentLoaded', () => {
  loadProducts();
});

// ===== Toast Notifications =====

function showToast(message, type = 'success') {
  const container = document.getElementById('toastContainer');

  const toast = document.createElement('div');
  toast.className = `toast toast-${type}`;
  toast.textContent = message;

  container.appendChild(toast);

  setTimeout(() => {
    toast.classList.add('hide');
    setTimeout(() => {
      toast.remove();
    }, 300);
  }, 3000);
}

// ===== Read Products =====

async function loadProducts() {
  const tbody = document.getElementById('productBody');

  tbody.innerHTML = '<tr><td colspan="6" style="text-align:center; color:#888;">Loading...</td></tr>';

  try {
    const { data: products, error } = await supabase
      .from('products')
      .select('*');

    if (error) throw error;

    renderProducts(products);
  } catch (error) {
    console.error('Products load karne me error:', error);
    tbody.innerHTML = '<tr><td colspan="6" style="text-align:center; color:#e74c3c;">Products load karne me error aaya. Server check karein.</td></tr>';
    showToast('Error: Server se connect nahi ho paya', 'error');
  }
}

// ===== Render Products =====

function renderProducts(products) {
  const tbody = document.getElementById('productBody');

  if (products.length === 0) {
    tbody.innerHTML = '<tr><td colspan="6" style="text-align:center; color:#888;">Koi product nahi hai</td></tr>';
    return;
  }

  const rows = products.map(product => `
    <tr>
      <td>${product.id}</td>
      <td>${product.name}</td>
      <td>Rs. ${product.price}</td>
      <td>${product.stock}</td>
      <td>${product.category}</td>
      <td>
        <button class="btn-edit" data-id="${product.id}" 
                data-name="${product.name}" data-price="${product.price}" 
                data-stock="${product.stock}" data-category="${product.category}">Edit</button>
        <button class="btn-delete" data-id="${product.id}">Delete</button>
      </td>
    </tr>
  `).join('');

  tbody.innerHTML = rows;
}

// ===== Modal Open/Close =====

document.getElementById('addProductBtn').addEventListener('click', () => {
  editingId = null;
  document.getElementById('modalTitle').textContent = 'Add Product';
  document.getElementById('productForm').reset();
  document.getElementById('productModal').classList.remove('hidden');
});

document.getElementById('cancelBtn').addEventListener('click', () => {
  document.getElementById('productModal').classList.add('hidden');
});

document.getElementById('productModal').addEventListener('click', (e) => {
  if (e.target === document.getElementById('productModal')) {
    document.getElementById('productModal').classList.add('hidden');
  }
});

// ===== Form Submit (Add/Edit) =====

let editingId = null;

document.getElementById('productForm').addEventListener('submit', async (e) => {
  e.preventDefault();

  const name = document.getElementById('productName').value.trim();
  const price = parseFloat(document.getElementById('productPrice').value);
  const stock = parseInt(document.getElementById('productStock').value);
  const category = document.getElementById('productCategory').value.trim();

  if (!name || !category) {
    showToast('Error: Name aur Category zaroori hain!', 'error');
    return;
  }

  if (isNaN(price) || price < 0) {
    showToast('Error: Price negative nahi ho sakta!', 'error');
    return;
  }

  if (isNaN(stock) || stock < 0) {
    showToast('Error: Stock negative nahi ho sakta!', 'error');
    return;
  }

  const productData = { name, price, stock, category };

  try {
    if (editingId !== null) {
      const { error } = await supabase
        .from('products')
        .update(productData)
        .eq('id', editingId);

      if (error) {
        showToast(`Error: ${error.message || 'Update me error aaya'}`, 'error');
        return;
      }
    } else {
      const { error } = await supabase
        .from('products')
        .insert([productData]);

      if (error) {
        showToast(`Error: ${error.message || 'Add karne me error aaya'}`, 'error');
        return;
      }
    }

    document.getElementById('productModal').classList.add('hidden');
    document.getElementById('productForm').reset();
    editingId = null;
    loadProducts();
    showToast('Product saved ho gaya!', 'success');
  } catch (error) {
    console.error('Error:', error);
    showToast('Error: Server se connect nahi ho paya', 'error');
  }
});

// ===== Edit & Delete Events =====

document.getElementById('productBody').addEventListener('click', (e) => {
  if (e.target.classList.contains('btn-edit')) {
    const id = e.target.dataset.id;
    const name = e.target.dataset.name;
    const price = e.target.dataset.price;
    const stock = e.target.dataset.stock;
    const category = e.target.dataset.category;

    editingId = parseInt(id);
    document.getElementById('modalTitle').textContent = 'Edit Product';

    document.getElementById('productName').value = name;
    document.getElementById('productPrice').value = price;
    document.getElementById('productStock').value = stock;
    document.getElementById('productCategory').value = category;

    document.getElementById('productModal').classList.remove('hidden');
  }

  if (e.target.classList.contains('btn-delete')) {
    const id = e.target.dataset.id;

    if (confirm('Kya aap yeh product delete karna chahte hain?')) {
      deleteProduct(id);
    }
  }
});

// ===== Delete Product =====

async function deleteProduct(id) {
  try {
    const { error } = await supabase
      .from('products')
      .delete()
      .eq('id', id);

    if (error) {
      showToast(`Error: ${error.message || 'Delete me error aaya'}`, 'error');
      return;
    }

    loadProducts();
    showToast('Product delete ho gaya!', 'success');
  } catch (error) {
    console.error('Error:', error);
    showToast('Error: Server se connect nahi ho paya', 'error');
  }
}
\end{lstlisting}

\end{document}